% UBT-AI-PROVENANCE-BEGIN
% schema: ubt-ai-provenance/v1
% tier: C_working
% ai_assistance: disclosed
% human_review: risk-based
% editorial_responsibility: Ing. David Jaroš
% policy: ../../../AI_PROVENANCE.md
% notice: Working material; exhaustive human review is not claimed.
% UBT-AI-PROVENANCE-END

\chapter{Continuing into complex time and Ward's identity}
\label{ch:ct-scheme}

\section{Why this chapter remains in the textbook}
Older alpha work used a ``CT scheme'' built around the analytic continuation
\(\tau=t+i\psi\).  The archived implementation is not a current source of
record, but the underlying mathematics---analytic continuation, gauge
identities, and the real-time limit---is standard and useful.  This chapter
therefore teaches those ingredients and then states exactly what they do and do
not establish for UBT.

\section{Complex time as a parameter}
For an analytic quantity $F(\tau)$,
\[
 \tau=t+i\psi
\]
allows one to move between oscillatory and damped descriptions.  For an energy
eigenmode, for example,
\[
 e^{-iE\tau/\hbar}=e^{-iEt/\hbar}e^{+E\psi/\hbar},
\]
with the sign of the damping determined by the convention for the imaginary
direction.  In the theta convention used later,
\[
 e^{-\pi(\psi-it)n^2}
 =e^{-\pi\psi n^2}e^{i\pi t n^2}.
\]
Nothing in this analytic identity by itself creates a new renormalisation
constant or fixes a physical coupling.

\section{The Ward--Takahashi identity: a standard argument}
In QED let $S(p)$ be the full fermion propagator and
$\Gamma^\mu(p+q,p)$ the full electromagnetic vertex.  Gauge invariance gives
the Ward--Takahashi identity
\[
 q_\mu\Gamma^\mu(p+q,p)
 =S^{-1}(p+q)-S^{-1}(p).
\]
Taking the limit $q\to0$ yields
\[
 \Gamma^\mu(p,p)=\frac{\partial S^{-1}(p)}{\partial p_\mu}.
\]
At the level of renormalisation constants this enforces
\[
 \boxed{Z_1=Z_2,}
\]
provided the regulator and subtraction prescription preserve gauge symmetry.
This is a standard gauge-theory result.  A complex-time continuation may use
it only after demonstrating that its contour/regulator preserves the hypotheses
required by the identity.

\section{The transversality of the photon's own energy}
The same gauge symmetry implies
\[
 q_\mu\Pi^{\mu\nu}(q)=0,
\]
so the vacuum polarisation has the form
\[
 \Pi^{\mu\nu}(q)
 =\left(q^2g^{\mu\nu}-q^\mu q^\nu\right)\Pi(q^2).
\]
(up to convention).  This is why charge renormalisation can be expressed in
terms of the transverse photon two-point function.  Again, this is a standard
field-theory constraint, not a UBT-specific theorem.

\section{Controlled UBT bridge}
A UBT complex-time renormalisation construction would need to establish, for
its actual action and contour:
\begin{enumerate}
 \item analyticity sufficient to continue between the chosen $\psi$ contour
       and the real-time slice;
 \item preservation of the relevant gauge/BRST identities;
 \item a regulator and counterterm scheme compatible with those identities;
 \item continuity to standard QED in the appropriate real-time limit;
 \item an explicit calculation of any finite CT-specific remainder rather than
       setting it by definition.
\end{enumerate}
Only after these steps could a CT-specific factor be promoted from a modelling
choice to a derived quantity.

\section{Why the archive text \(\mathcal R_{\rm UBT}=1\) is not downloaded}
The historical alpha branch stated a two-loop baseline
\(\mathcal R_{\rm UBT}=1\) under assumptions labelled A1--A3 and then called the
resulting alpha baseline ``fit-free''.  Current Tier-A status documents no
longer allow that wording to stand as a completed alpha derivation: the
coefficient bridge G137-B is open and the effective multiplicity has an audit
gap.  Therefore the live textbook keeps the Ward-identity mathematics but does
not reuse the archived theorem as a current source of record.

This is a useful example of a general rule: a correct conditional field-theory
identity can survive even when a larger phenomenological interpretation built
on top of it is downgraded.

\section{Link to Theta Chapter}
The reduced theta kernel
\[
 S(t,\psi)=\sum_n a_n e^{-\pi\psi n^2}e^{i\pi t n^2}
\]
provides a cleaner controlled laboratory for complex time than the old alpha
renormalisation route.  The same parameter $\psi$ acts as a heat-time variable,
while $t$ carries the oscillatory phase.  Fourier, Laplace, Mellin, and
Poisson/Jacobi transforms can then be derived exactly before any physical UBT
claim is attached to them.
